%% legal-document.tex
%% Template generico per documenti legali — PoliNetwork APS
%% 
%% Adatto a: statuto, regolamento interno, delibere, atti formali
%%
%% Compila con: xelatex legal-document.tex
%%
\documentclass[legal]{polinetwork}

%% ====================================================================
%% DATI DOCUMENTO
%% ====================================================================

\pnTitle{[Titolo del Documento]}
\pnSubtitle{[Sottotitolo opzionale]}
\pnDate{1 aprile 2026}
\pnVersion{1.0}
\pnSubject{[Oggetto del documento]}
\pnProtocol{PN-2026/001}

%% Per documenti in inglese, decommentare:
% \pnSetEnglish

%% ====================================================================

\begin{document}
\thispagestyle{firstpage}

\pnMakeHeader
\pnMakeMetadata

%% ====================================================================
%% PREMESSA / PREAMBOLO (opzionale)
%% ====================================================================

% Decommentare se serve una premessa
% \pnSection{Premessa}
% Testo della premessa...

%% ====================================================================
%% ARTICOLI
%% ====================================================================

\pnArticle{[Titolo articolo 1]}

\pnClause{%
  [Testo della prima clausola dell'articolo 1.]
}

\pnClause{%
  [Testo della seconda clausola dell'articolo 1.]
}

\pnArticle{[Titolo articolo 2]}

\pnClause{%
  [Testo della prima clausola dell'articolo 2.]
}

% Puoi anche usare elenchi all'interno delle clausole:
\pnClause{%
  [Introduzione alla clausola con elenco]:
  \begin{itemize}
    \item primo punto;
    \item secondo punto;
    \item terzo punto.
  \end{itemize}
}

\pnArticle{[Titolo articolo 3]}

\pnClause{%
  [Testo...]
}

%% ====================================================================
%% CHIUSURA E FIRMA
%% ====================================================================

\pnClosing

\vfill
\pnContactBlock

\end{document}
